\phantomsection\addcontentsline{toc}{section}{Additive Manufacturing}
\ECchapter{11}{Additive Manufacturing}{Can 3D printing transform industrial production?}{ECOrange}

\ECObjectives{
\begin{itemize}
\item Explain the complete digital chain from CAD model to qualified part.
\item Compare additive manufacturing with machining and casting using engineering criteria.
\item Justify design, orientation, post-processing and inspection choices for a printed component.
\end{itemize}}

\begin{multicols}{2}
\begin{engineeringnews}
Additive manufacturing is now used for rocket engines, aircraft components, medical implants and
industrial tooling. Its real value is not simply the ability to print a shape, but the possibility of
redesigning a function, reducing assembly operations and producing geometries that conventional
processes cannot manufacture economically.
\end{engineeringnews}

\begin{ecarticle}
Additive manufacturing creates a component layer by layer from a digital model. The process begins
with a three-dimensional CAD file. Slicing software divides the geometry into thin horizontal
sections and generates machine instructions. Depending on the technology, a nozzle deposits a
molten polymer, a laser melts metal powder, or light solidifies a liquid resin.

This manufacturing route changes the designer's freedom. Internal cooling channels, lattice
structures and organic shapes can be integrated into a single component. Topology optimisation can
remove material from low-stress zones while preserving principal load paths. In aerospace
engineering, this can reduce mass and fuel consumption. In medicine, a porous implant can be adapted
to a patient and encourage bone growth.

The process also has important constraints. Overhanging zones may require temporary supports. These
supports consume material, conduct heat and must be removed after printing. Thermal gradients can
create residual stress, warping or cracks, especially in metal parts. Build orientation therefore
affects support volume, surface quality, printing time, dimensional accuracy and mechanical
properties.

A printed component is rarely ready for immediate use. Powder removal, heat treatment, support
removal, machining and inspection may be required. Non-destructive testing can detect internal
porosity, while dimensional measurement verifies that the geometry remains within tolerance. Since
mechanical properties may vary with build direction, qualification is essential for safety-critical
parts.

Additive manufacturing is particularly effective for prototypes, customised products, spare parts
and low-volume high-value components. It is less competitive for millions of simple identical parts.
The best industrial solution is often hybrid: printing creates the complex near-net shape, then
conventional machining produces accurate interfaces and smooth functional surfaces.
\end{ecarticle}

\begin{tcolorbox}[title=\sffamily\bfseries Engineering Insight,colback=yellow!10,colframe=ECOrange]
A successful printed part is designed together with its manufacturing plan. Geometry, orientation,
supports, scan strategy, heat treatment, machining references and inspection method form one
integrated engineering decision.
\end{tcolorbox}

\begin{engineeringfocus}
\textbf{Digital manufacturing chain}

\begin{center}
\begin{tikzpicture}[font=\scriptsize,>=Latex,node distance=3.2mm,scale=.92,transform shape]
\tikzset{am/.style={draw=ECOrange,fill=ECOrange!9,rounded corners,align=center,
minimum width=22mm,minimum height=10mm}}
\node[am] (cad) {CAD model};
\node[am,right=of cad] (slice) {slicing and\\supports};
\node[am,right=of slice] (build) {layer-by-layer\\build};
\node[am,below=of build] (post) {post-processing};
\node[am,left=of post] (inspect) {inspection and\\qualification};
\node[draw=ECGreen,fill=ECGreen!10,rounded corners,align=center,left=of inspect,
minimum width=22mm,minimum height=10mm] (part) {finished part};
\draw[->,thick] (cad) -- (slice);
\draw[->,thick] (slice) -- (build);
\draw[->,thick] (build) -- (post);
\draw[->,thick] (post) -- (inspect);
\draw[->,thick] (inspect) -- (part);
\draw[->,dashed,ECBlue] (inspect.north) |- (cad.south);
\end{tikzpicture}
\end{center}

Inspection may lead to a design or parameter change. The chain is therefore iterative rather than
strictly linear.
\end{engineeringfocus}

\begin{keyfigures}
\begin{tabularx}{\linewidth}{@{}Xr@{}}
Typical polymer layer & 0.10--0.30 mm\\
Typical metal layer & 0.02--0.08 mm\\
Main digital input & 3D CAD model\\
Frequent post-processes & heat treatment, machining\\
Best production range & low volume / high value
\end{tabularx}
\end{keyfigures}

\begin{tcolorbox}[title=\sffamily\bfseries Comparing manufacturing routes,colback=white,colframe=ECOrange]
\begin{center}
\begin{tikzpicture}
\begin{axis}[
width=.94\linewidth,height=45mm,
ybar,bar width=7pt,
ymin=0,ymax=105,
ylabel={Relative score},
symbolic x coords={Machining,Casting,Additive manufacturing},
xtick=data,x tick label style={font=\tiny,rotate=18,anchor=east},
yticklabel style={font=\scriptsize},
legend style={font=\fontsize{4.5}{5}\selectfont,at={(0.5,1.02)},anchor=south,legend columns=3},
grid=major]
\addplot table[x=process,y=time,col sep=comma]{data/EC11/process_comparison.csv};
\addplot table[x=process,y=material,col sep=comma]{data/EC11/process_comparison.csv};
\addplot table[x=process,y=complexity,col sep=comma]{data/EC11/process_comparison.csv};
\legend{Speed,Material efficiency,Geometric freedom}
\end{axis}
\end{tikzpicture}
\end{center}
\footnotesize Illustrative comparison. The most suitable process depends on quantity, material,
tolerance, geometry, qualification effort and total lifecycle cost.
\end{tcolorbox}

\begin{tcolorbox}[title=\sffamily\bfseries Designing a lightweight bracket,colback=white,colframe=ECBlue]
\begin{center}
\begin{tikzpicture}[scale=.82]
\draw[very thick,ECBlue,rounded corners=3pt] (0,0) rectangle (4.8,2.6);
\fill[white] (.9,.8) circle (.38); \draw[thick] (.9,.8) circle (.38);
\fill[white] (3.9,.8) circle (.38); \draw[thick] (3.9,.8) circle (.38);
\fill[white] (2.4,1.9) circle (.42); \draw[thick] (2.4,1.9) circle (.42);
\foreach \x/\y in {1.55/.75,2.4/.75,3.25/.75,1.75/1.45,3.05/1.45}{
  \draw[ECOrange,thick] (\x,\y) -- ++(.42,.34) -- ++(-.42,.34) -- ++(-.42,-.34) -- cycle;}
\draw[->,thick,ECRed] (2.4,3.35) -- (2.4,2.45) node[midway,right]{load};
\node[font=\scriptsize,align=center] at (2.4,-.45) {material kept around interfaces\\and principal load paths};
\end{tikzpicture}
\end{center}
Topology optimisation must be followed by checks on minimum wall thickness, support removal,
powder evacuation and machinable reference surfaces.
\end{tcolorbox}

\begin{vocabulary}
\begin{tabularx}{\linewidth}{@{}lX@{}}
additive manufacturing & fabrication additive\\
slicing & découpage en couches\\
powder bed & lit de poudre\\
nozzle & buse\\
support structure & structure de support\\
build orientation & orientation de fabrication\\
warping & déformation / gauchissement\\
porosity & porosité\\
post-processing & post-traitement\\
topology optimisation & optimisation topologique
\end{tabularx}
\end{vocabulary}

\begin{questions}
\item What information does slicing software generate?
\item Why does build orientation influence support volume and mechanical performance?
\item Give three examples of post-processing or inspection operations.
\item Why can mechanical properties be anisotropic?
\item Use the comparison chart to identify one advantage and one limitation of each process.
\item Why may part consolidation improve reliability but reduce repairability?
\end{questions}

\begin{applicationexercise}
A metal part is 120 mm high and is printed with a layer thickness of 0.04 mm. Each layer takes
18 seconds to scan and recoating adds no extra time in this simplified model.
\begin{enumerate}
\item Calculate the number of layers.
\item Estimate the ideal build time in hours.
\item Add 25\% for calibration and interruptions. Estimate the total production time.
\item Explain why the real industrial lead time also includes preparation, cooling and inspection.
\end{enumerate}
\end{applicationexercise}

\begin{engineeringchallenge}
Redesign an aluminium bracket for metal additive manufacturing. Define loads, interfaces and
acceptance criteria; propose a lighter geometry; select a build orientation; estimate support and
post-processing needs; identify inspection methods; and compare the total production route with
machining from solid. Conclude with a justified recommendation.
\end{engineeringchallenge}

\begin{speaking}
Will additive manufacturing replace conventional factories? Give a balanced four-minute answer with
examples of suitable and unsuitable products, and distinguish technical possibility from economic
relevance.
\end{speaking}
\end{multicols}

\subsection*{Sources}
\ECsource{NIST -- Additive Manufacturing}{https://www.nist.gov/additive-manufacturing}
\ECsource{NASA -- Additive Manufacturing}{https://www.nasa.gov/centers-and-facilities/marshall/additive-manufacturing/}
